% Source: https://artofproblemsolving.com/community/c3h2673524_question
% Stable identity: aops-topic-2673524-q01
$18^n=2^n3^{2n}$. By Legendre's formula the exponents of $2$ and $3$ in $200!$ are
\[\left\lfloor\dfrac{200}{2}\right\rfloor+\left\lfloor\dfrac{200}{4}\right\rfloor+\cdots=100+50+25+12+6+3+1=197,\]
\[\left\lfloor\dfrac{200}{3}\right\rfloor+\left\lfloor\dfrac{200}{9}\right\rfloor+\left\lfloor\dfrac{200}{27}\right\rfloor+\left\lfloor\dfrac{200}{81}\right\rfloor=66+22+7+2=97.\]
We need $n\le197$ and $2n\le97$, so $n=\boxed{48}$.

$\left\lfloor\dfrac{200}{18}\right\rfloor$ counts only the multiples of $18$ up to $200$. But factors of $2$ and $3$ coming from different numbers also combine into factors of $18$ (for example, the factor $2$ of $2$ and the factor $9$ of $9$, although neither $2$ nor $9$ is a multiple of $18$), so $18$ divides $200!$ far more often than there are multiples of $18$. Only for a prime $p$ does counting multiples (with powers) give the exponent.
